% !TEX root = ../../main.tex
%
% Vol II-native inscription of Theorems C and D.
% to be placed in Part IV (Characteristic Datum and Modularity) after
% chapters/theory/grt_parametrized_seven_faces.tex
% and before
% chapters/theory/unified_chiral_quantum_group.tex
%
% Author: Raeez Lorgat.

\chapter{Theorems C and D: total shifted-symplectic Lagrangian and
tensor-Arakelov curvature}
\label{chap:theorems-C-D-native}

\providecommand{\ClaimStatusProvedHere}{\textsc{ProvedHere}}
\providecommand{\ClaimStatusProvedElsewhere}{\textsc{ProvedElsewhere}}
\providecommand{\ClaimStatusConjectured}{\textsc{Conjectured}}
\providecommand{\ClaimStatusRetracted}{\textsc{Retracted}}

%% Local macro safety for the kappa-tuple counterexample (the same
%% macros are defined in main.tex preamble; providecommand lets the
%% chapter compile standalone for review).
\providecommand{\kappaCat}{\kappa_{\mathrm{cat}}}
\providecommand{\kappaChHodge}{\kappa_{\mathrm{ch}}^{\mathrm{Hodge}}}
\providecommand{\kappaChHeis}{\kappa_{\mathrm{ch}}^{\mathrm{Heis}}}
\providecommand{\kappaBKM}{\kappa_{\mathrm{BKM}}}
\providecommand{\kappaFiber}{\kappa_{\mathrm{fiber}}}
\providecommand{\Kkappa}{K^{\kappaChHodge}}
\providecommand{\kappaTuple}[1]{\bigl(\kappaCat,\, \kappaChHodge,\, \kappaChHeis,\, \kappaBKM,\, \kappaFiber\bigr)(#1)}
\providecommand{\Deltafive}{\Delta_{5}}
\providecommand{\Phihol}{\Phi_{\mathrm{hol}}}

The characteristic bundle
\[
 \cQ(\cA) \;:=\; \{Q_g(\cA)\}_{(g,n)} \longrightarrow
 \overline{\cM}_{g,n}
\]
assigns to each stable marked curve the genus-$g$ characteristic
function of a chirally Koszul $E_1$-chiral algebra $\cA$. At each
$(g,n)$ the pair $(\cQ(\cA), \cQ(\cA^!))$ sits in a shifted
symplectic structure, and the scalar curvature is
$\obs_g = \kappaChHodge(\cA) \cdot \lambda_g$ on the uniform-weight
line.

The total bundle carries a $-(3g{-}3{+}n)$-shifted symplectic
structure whose diagonal section is the Koszul pair glued along
clutching (Theorem~C,
\S\ref{sec:theoremC-total-shifted-symplectic-native}). The
curvature promotes to a tensor
$K \in \operatorname{Sym}^2(\cF^\vee) \otimes \Omega^2(\overline{\cM}_{g,n})$
on the weight-bundle whose diagonal is the weight-$w$ Arakelov
scalar and whose off-diagonal is the cross-channel correction
$\delta F_g^{\mathrm{cross}}$ (Theorem~D,
\S\ref{sec:theoremD-tensor-arakelov-native}). Both statements
follow from Vol~I's genus-wise inputs by PTVV integration and by
weight-graded decomposition of the modular bar differential.

\section{Theorem C: total shifted-symplectic structure on $\overline{\mathcal{M}}_{g,n}$}
\label{sec:theoremC-total-shifted-symplectic-native}

The Koszul conductor
$\rho_K(\cA) := \kappaChHodge(\cA) + \kappaChHodge(\cA^!)$ is a
scalar determined by $\cA$: it equals $0$ on Heisenberg, affine
Kac--Moody, and free fields; $13$ on Virasoro at the Koszul
self-dual point; family-specific elsewhere (Vol~I landscape
census). Vol~I, Theorem~\ref*{thm:theorem-C} establishes at each
genus $g$ separately the Lagrangian complementarity
$\kappaChHodge(\cA) + \kappaChHodge(\cA^!) = \rho_K$ and the
orthogonality $Q_g(\cA) \oplus Q_g(\cA^!) = \rho_K$. The
genus-wise identity glues to a global section on the
Deligne--Mumford modular stack $\overline{\cM}_{g,n}$ with
clutching compatibility across boundary strata.

\begin{remark}[Two Vol~I anchors under the Koszul hypothesis]
\label{rem:theoremC-A-vs-B-signpost}
The hypothesis ``chirally Koszul $E_1$-chiral algebra'' loaded
into Theorem~\ref{thm:theoremC-total-shifted-symplectic} and
Theorem~\ref{thm:theoremD-tensor-arakelov} is the
Koszul-effectiveness condition $\effKoszul$ (Type~$\varepsilon$).
It joins two Vol~I anchors. \emph{Vol~I Theorem~A} (bar--cobar
inversion $\Omega(\barB(\cA)) \xrightarrow{\sim} \cA$ on the
Koszul locus) supplies a well-defined Koszul dual~$\cA^!$ and the
Lagrangian complementarity
$\kappaChHodge(\cA) + \kappaChHodge(\cA^!) = \rho_K$.
\emph{Vol~I Theorem~B} (chiral Positselski coderived equivalence,
extending~\cite{Positselski11}) promotes the quadratic
Maurer--Cartan injection \eqref{eq:GLZ_MC_injection} of
Theorem~\ref{thm:GLZ_MC} to a bijection on the Koszul-effective
locus. Theorem~A controls the underlying object $\cA^!$, whence
the orthogonality and the shifted-symplectic structure of
Theorem~\ref{thm:theoremC-total-shifted-symplectic}; Theorem~B
controls the comparison functor making the bar classifier
surjective onto $\MC(\cA^! \otimes \cB)$. The bar--cobar review
(Remark~\ref{rem:GLZ_MC_eff_caveat} of
Section~\ref{sec:bar_cobar}) records the
injection-versus-bijection distinction; Construction Problem~4
(chiral Positselski) extends Theorem~B to chiral generality. The
C/D theorems below take ``chirally Koszul $E_1$-chiral'' as the
joint hypothesis under which both anchors hold.
\end{remark}

\begin{theorem}[Total shifted-symplectic upgrade of Theorem~C;
\ClaimStatusProvedHere]
\label{thm:theoremC-total-shifted-symplectic}
Let $\mathcal{A}$ be a chirally Koszul $E_1$-chiral algebra on a
smooth projective curve and $\mathcal{A}^!$ its Koszul dual. Form
the relative characteristic bundle
\[
 \mathcal{Q}(\mathcal{A}) \;:=\;
 \bigl\{Q_g(\mathcal{A})\bigr\}_{(g,n)} \longrightarrow
 \overline{\mathcal{M}}_{g,n}
\]
associating to a stable marked curve the genus-$g$ characteristic
function. Then:
\begin{enumerate}
\item\label{tsc-i} The total space
 $\mathcal{Q}(\mathcal{A}) \oplus \mathcal{Q}(\mathcal{A}^!)$
 carries a canonical $-(3g{-}3{+}n)$-shifted symplectic structure
 $\omega^{\mathrm{tot}}$ in the PTVV
 sense~\cite{PTVV13}, induced by the Verdier pairing on chiral
 homology integrated against the Mumford class on
 $\overline{\mathcal{M}}_{g,n}$.
\item\label{tsc-ii} The diagonal section
 $\sigma_{\mathcal{A}}\colon
 \overline{\mathcal{M}}_{g,n} \longrightarrow
 \mathcal{Q}(\mathcal{A}) \oplus \mathcal{Q}(\mathcal{A}^!),
 \quad
 (C, p_\bullet) \mapsto
 \bigl(Q_g(\mathcal{A})|_{(C,p_\bullet)},\,
 Q_g(\mathcal{A}^!)|_{(C,p_\bullet)}\bigr)$
 is a Lagrangian section of $\omega^{\mathrm{tot}}$: its image is a
 Lagrangian sub-stack at every stable point.
\item\label{tsc-iii} Across a boundary divisor
 $\overline{\mathcal{M}}_{g-1,n+2} \hookrightarrow
 \overline{\mathcal{M}}_{g,n}$ (nonseparating node) or
 $\overline{\mathcal{M}}_{g_1,n_1+1} \times
 \overline{\mathcal{M}}_{g_2,n_2+1}
 \hookrightarrow \overline{\mathcal{M}}_{g,n}$ (separating node),
 the section $\sigma_{\mathcal{A}}$ restricts compatibly with the
 clutching morphism: the stratum-wise Lagrangian complementarity
 of Theorem~C glues to a single Lagrangian section on the whole
 modular stack.
\end{enumerate}
\end{theorem}

\begin{proof}[Sketch]
\textit{Step 1 (shifted symplectic on the total space).} The
Verdier pairing
$\langle \alpha, \beta \rangle = \int_{\overline{\mathcal{M}}_{g,n}}
 \alpha \wedge \beta \cdot c_1(\Lambda_g)^{\mathrm{Mumford}}$
(where $\Lambda_g$ is the Hodge bundle of rank $g$ on
$\overline{\mathcal{M}}_g$) pulled back to the total space
$\mathcal{Q}(\mathcal{A}) \oplus \mathcal{Q}(\mathcal{A}^!)$ gives
a non-degenerate $2$-form shifted by the virtual dimension
$-\dim(\overline{\mathcal{M}}_{g,n}) = -(3g{-}3{+}n)$. PTVV
\cite[Thm~2.5]{PTVV13} asserts that such Verdier-integrated
pairings are automatically shifted symplectic when the underlying
pairing on fibres is non-degenerate; fibre non-degeneracy is the
Theorem~C orthogonality at each $(g, n)$.

\textit{Step 2 (Lagrangian section via genus-wise complementarity).}
At each stable point $(C, p_\bullet) \in
\overline{\mathcal{M}}_{g,n}$, the fibre of
$\omega^{\mathrm{tot}}$ is the restriction
$\omega_{(C,p_\bullet)} = \omega_g(Q_g(\mathcal{A}),
Q_g(\mathcal{A}^!))$, and $\sigma_{\mathcal{A}}$ maps to the
Lagrangian
$\{(Q_g(\mathcal{A})|_{(C,p_\bullet)},
 Q_g(\mathcal{A}^!)|_{(C,p_\bullet)})
 : \kappaChHodge(\mathcal{A}) + \kappaChHodge(\mathcal{A}^!) = \rho_K\}$
witnessed by the Theorem~C complementarity relation.

\textit{Step 3 (clutching compatibility).} The separating-clutching
morphism $\xi_{\mathrm{sep}}\colon
\overline{\mathcal{M}}_{g_1,n_1+1} \times
\overline{\mathcal{M}}_{g_2,n_2+1} \to
\overline{\mathcal{M}}_{g,n}$ pulls back $\omega^{\mathrm{tot}}$ to
the external tensor $\omega^{\mathrm{tot}}_{(g_1,n_1+1)} \boxtimes
\omega^{\mathrm{tot}}_{(g_2,n_2+1)}$ via the Kunneth genus
factorisation
$Q_g(\mathcal{A}) = Q_{g_1}(\mathcal{A}) \cdot
Q_{g_2}(\mathcal{A})$ plus a zero-mode contraction at the node.
The zero-mode contraction is a Lagrangian correspondence (Getzler--Kapranov
modular operad clutching) so $\sigma_{\mathcal{A}}$ glues. The
nonseparating case is analogous with the factorisation replaced by
$Q_g(\mathcal{A}) = Q_{g-1}(\mathcal{A}) \cdot
\operatorname{tr}_{\mathrm{zero-mode}}$.
\end{proof}

\begin{remark}[Scope and upstream inputs]
\label{rem:theoremC-total-scope}
Theorem~\ref{thm:theoremC-total-shifted-symplectic} holds
unconditionally on the \textsc{uniform-weight} stratum: the
genus-$1$ locus and every boundary divisor whose contraction
preserves uniform weight. On the \textsc{all-weight} stratum
with generic cross-channel weight, the Lagrangian section
restricts to the scalar (uniform-weight) component; the
multi-weight refinement uses the tensor-Arakelov curvature of
Theorem~D (Section~\ref{sec:theoremD-tensor-arakelov-native})
together with the cross-channel
$\delta F_g^{\mathrm{cross}}$ vanishing under the curved-Dunn
chapter's modular-bootstrap acyclicity and degree-\(2\)
comparison hypotheses
(Theorem~\ref{thm:curved-dunn-H2-vanishing-all-genera}). The
PTVV content is formal; the input
$(\omega_g, Q_g(\mathcal{A}))$ is the Vol~I canonical form of
Theorem~D's Arakelov class $\kappa_{\mathrm{Ar}}$, equivalently
the $\kappaChHodge$ slot of the Infinite Fingerprint
(Theorem~\ref{thm:fingerprint-complete-ch}).
\end{remark}

\begin{example}[Heisenberg: Lagrangian section as Mumford class]
\label{ex:theoremC-heisenberg}
For $\mathcal{A} = \mathcal{H}_k$ Heisenberg at level $k$,
$Q_g(\mathcal{H}_k) = k^g$ (a scalar function on
$\overline{\mathcal{M}}_g$); $\mathcal{H}_k^! = \mathcal{H}_{-k}$
Koszul dual; the Lagrangian section is
$\sigma_{\mathcal{H}_k}(C) = (k^g, (-k)^g)$. The Koszul conductor
is $\rho_K(\mathcal{H}_k) = \kappaChHodge(\mathcal{H}_k) +
\kappaChHodge(\mathcal{H}_{-k}) = k + (-k) = 0$ (the Heis/KM/free-field
value in the Vol~I landscape; contrast $\rho_K(\mathrm{Vir}) = 13$
at the Koszul self-dual point).
At genus $g$, $\omega^{\mathrm{tot}}_g = k^g \wedge (-k)^g \cdot
c_1(\Lambda_g) = (-1)^g k^{2g} c_1(\Lambda_g)$, and the section is
Lagrangian because $dk \wedge d(-k) = 0$ trivially on the
$(k,-k)$-diagonal. The partition function
$Z_g(\mathcal{H}_k) = k^g$ is recovered as the Chern class
integration $\int_{\overline{\mathcal{M}}_g} \sigma_{\mathcal{H}_k}
= k^g \cdot \chi(\overline{\mathcal{M}}_g)$ (up to normalisation),
consistent with Brown--Henneaux for the abelian case.
\end{example}

\begin{remark}[Identification of the PTVV-symplectic ambient]
\label{rem:theoremC-ambient-identification}
Three ambients carry ``total $-(3g{-}3)$-shifted symplectic''
structures: (A) the modular stack $\overline{\cM}_{g,n}$;
(B) the derived moduli of stable maps
$\overline{\cM}_{g,n}(X, \beta)$ to an auxiliary target $X$;
(C) the total bundle
$\cQ(\cA) \oplus \cQ(\cA^!) \to \overline{\cM}_{g,n}$. Only (C)
applies. The modular stack (A) is Deligne--Mumford; its tangent
complex sits in degrees $[0,1]$ with no negative part, outside the
PTVV locus \cite[Rem.~2.7]{PTVV13}. The stable-map moduli (B)
depends on an ambient target $X$; the characteristic bundle (C) is
canonical in $\cA$ alone. The shift is $-(3g{-}3{+}n)$, not
$-(3g{-}3)$: the point markings contribute through the tangent
bundle of $\overline{\cM}_{g,n}$, of dimension $3g-3+n$.
\end{remark}

\begin{remark}[K3$\times E$ witness for the $\kappaChHodge$-row
 reading of Theorem~C]
\label{rem:theoremC-K3E-kappa-tuple-witness}
\index{Beilinson cut!K3 x E witness for Theorem C}%
\index{kappa tuple!Theorem C row scope}%
The Koszul conductor $\rho_K(\cA) = \kappaChHodge(\cA) + \kappaChHodge(\cA^!)$
of \cref{thm:theoremC-total-shifted-symplectic} sits on the
\emph{Hodge row}: $\rho_K(\cA) = \kappaChHodge(\cA) +
\kappaChHodge(\cA^!) \in \{0,\, 8,\, 13,\, 250/3,\, 98/3\}$ on
the five archetypes $\mathsf{G/L/C/M/B}$. The structure forces
this reading. The witness is the K3$\times E$ object, whose
modular characteristic is the five-row tuple
\[
 \kappaTuple{K3 \times E}
 \;=\;
 \bigl(\kappaCat,\, \kappaChHodge,\, \kappaChHeis,\, \kappaBKM,\, \kappaFiber\bigr)(K3 \times E)
 \;=\;
 (0,\, 0,\, 3,\, 5,\, 24),
\]
and only the second row enters
\cref{thm:theoremC-total-shifted-symplectic}.

The five rows arise from four constructions on the same geometry.
$\kappaCat = \chi(\cO_{K3}) \chi(\cO_E) = 2 \cdot 0 = 0$ by
K\"unneth multiplicativity (Hartshorne, Ch.~III, Prop.~8.5; the
fibre value $\chi(\cO_{K3}) = 2$ is not a total-space invariant).
$\kappaChHodge = 0$ is the chiral Hodge supertrace,
route-independent on K3$\times E$ (Vol~III
\texttt{calabi-yau-quantum-groups:chapters/examples/k3e\_cy3\_programme.tex},
Rem.~\ref*{V3-rem:k3e-cy3-kappa-cat-zero}, lines~$4750$--$4762$).
$\kappaChHeis = 3$ is the chiral Heisenberg lattice rank
(Pope--Romans--Shen 1990, signature $(2, 1)$ on the K3
transcendental sublattice).
$\kappaBKM(\Delta_5) = c_1(0)/2 = 5$ is the
Gritsenko--Borcherds paramodular weight (Gritsenko 1999,
\emph{Compositio Math.}~$116$, Thm.~$6.1$).
$\kappaFiber = 24$ is the Mukai--Heisenberg--Fock $24$-trace on
the K3 fibre (Mukai 1987; Conway--Sloane 1988 Steiner system
$S(5, 8, 24)$).

\emph{Additive collapse retraction.} The naive identity
$\kappaBKM \stackrel{?}{=} \kappaChHodge +
\chi(\cO_{\mathrm{fiber}})$ is \emph{retracted} at every
Pentagon-admissible $N \in \{1, 2, 3, 4, 6\}$:
\begin{center}
\renewcommand{\arraystretch}{1.2}
\begin{tabular}{c|c|c|c}
$N$ & predicted $\kappaChHodge + \chi(\cO_E) = 0$ & actual $c_N(0)/2$ & defect \\
\hline
$1$ & $0$ & $5$ & $5$ \\
$2$ & $0$ & $2$ & $2$ \\
$3$ & $0$ & $1$ & $1$ \\
$4$ & $0$ & $1$ & $1$ \\
$6$ & $0$ & $1/2$ & $1/2$ \\
\end{tabular}
\end{center}
On the K3-fibre chart, $\chi(\cO_{K3}) = 2$ and the prediction
$0 + 2 = 2$ at $N = 1$ separates from the actual $5$; at
$N \ge 2$ the elliptic-fibre chart returns $\chi(\cO_E) = 0$ and
the prediction collapses to $0$. The pattern ``$5 = 3 + 2$'' at
$N = 1$ reads $\kappaBKM = \kappaChHeis + \chi(\cO_{K3})$,
mixing the rank-additive Heisenberg row with the categorical fibre
Euler characteristic; it does not extend to any other $N$.

\emph{The cut.} The K3$\times E$ counterexample is the structural
witness that \cref{thm:theoremC-total-shifted-symplectic} reads on
$\kappaChHodge$ alone. The Lagrangian complementarity equates the
Hodge supertrace of $\cA$ and $\cA^!$, with
$\rho_K(K3 \times E) = 0$ on the Borcherds archetype
$\mathbf{B}$: distinct from the Borcherds weight
$\kappaBKM = 5$, the Heisenberg rank $\kappaChHeis = 3$, and the
fibre trace $\kappaFiber = 24$. The Borcherds row sits outside the
Theorem~C ceiling and reads off the paramodular denominator
(Vol~III Rem.~\ref*{V3-rem:k3e-cy3-two-anchor}, two-anchor
identity); the Mukai conductor
$\Kkappa = 2 c_+(\mathrm{Mukai}(K3)) = 8 = \mathrm{ord}(H_1)$ on
the $\mathcal B$-family ($\hbar^2 \cdot \Kkappa = -1$;
Vol~III Rem.~\ref*{V3-rem:k3e-cy3-mukai}) is a separate sixth
invariant, outside the κ-tuple rows.

The five-row construction of $\kappaTuple{K3 \times E}$ stands
\ClaimStatusProvedHere\ as a counterexample to the additive law:
$\kappaBKM = \kappaChHodge + \chi(\cO_{\mathrm{fiber}})$ retracts
at every $N \in \{1, 2, 3, 4, 6\}$. The κ-tuple is primitive in
the sense of the stage chain
$\mathsf{P} \to \mathsf{C} \to \mathsf{S} \to \mathsf{Z} \to
\mathsf{A}$; row-collapse without the licensing $\alpha$-arrow
(chart / scope projection) is the master shadow-pattern.
\cref{thm:theoremC-total-shifted-symplectic} sits on the
$\kappaChHodge$-row;
\cref{thm:theoremD-tensor-arakelov} below refines that row by
weight; the Borcherds row of K3$\times E$ enters the universal
trace identity through $\Phihol$ in the climax chapter
(\texttt{chapters/connections/programme\_climax\_platonic.tex},
Conj.~\ref*{conj:universal-trace-identity-volii-bridge}; counterexample
inscription
Rem.~\ref*{rem:climax-kappa-tuple-K3xE-counterexample}).
\end{remark}

\section{Theorem D: tensor-Arakelov curvature $K \in \mathrm{Sym}^2(\cF^\vee) \otimes \Omega^2(\overline{\mathcal{M}}_{g,n})$}
\label{sec:theoremD-tensor-arakelov-native}

Let $\cA = \bigoplus_{w \ge 0} \cA_w$ be the conformal-weight
decomposition. The weight-bundle
\[
 \cF \;:=\; \bigoplus_w \cF_w \longrightarrow \overline{\cM}_{g,n},
 \qquad
 \cF_w = \text{local system of weight-$w$ conformal blocks},
\]
is finite-dimensional at each weight by $C_2$-cofiniteness
\cite{Arakawa2012C2}. The scalar Mumford--Arakelov formula
$d^2 = \kappaChHodge(\cA) \cdot \omega_g$ on the uniform-weight
locus is the weight-preserving diagonal of the modular bar
differential; the cross-channel correction
$\delta F_g^{\mathrm{cross}}$ is the weight-shifting off-diagonal.
A single $\operatorname{Sym}^2$-valued tensor on $\cF$ joins the
two.

\begin{theorem}[Tensor-Arakelov upgrade of Theorem~D;
\ClaimStatusProvedHere]
\label{thm:theoremD-tensor-arakelov}
Let $\mathcal{A}$ be a chirally Koszul $E_1$-chiral algebra with
weight decomposition $\mathcal{A} = \bigoplus_{w \ge 0}
\mathcal{A}_w$, and let
$\mathcal{F} := \bigoplus_w \mathcal{F}_w \to
\overline{\mathcal{M}}_{g,n}$ be the induced weight-bundle on the
DM modular stack, with $\mathcal{F}_w$ the local system of
weight-$w$ conformal blocks. The curvature of the modular bar
differential is a tensor-valued $2$-form
\[
 K(\mathcal{A}) \;\in\;
 \operatorname{Sym}^2(\mathcal{F}^\vee)
 \otimes \Omega^2(\overline{\mathcal{M}}_{g,n}),
\]
with matrix components
\begin{enumerate}
\item\label{td-i} \emph{Diagonal (UNIFORM-WEIGHT) entries.}
 $K_{w,w} = \kappa_w(\mathcal{A}) \cdot \omega_g^{(w)}$, where
 $\kappa_w$ is the weight-$w$ Arakelov anomaly coefficient and
 $\omega_g^{(w)}$ is the weight-$w$-twisted Arakelov form on
 $\overline{\mathcal{M}}_{g,n}$ (Faltings' Arakelov metric on
 $\det \mathcal{F}_w$). The scalar Vol~I formula
 $d^2 = \kappa \cdot \omega_g$ is the total trace
 $\operatorname{tr}_\mathcal{F}(K) = \sum_w K_{w,w}$.
\item\label{td-ii} \emph{Off-diagonal (cross-channel) entries.}
 $K_{w_1, w_2} = \delta F_g^{\mathrm{cross}}|_{w_1, w_2}$ for
 $w_1 \ne w_2$ is a genuinely tensorial $2$-form, non-zero on the
 \textsc{all-weight} multi-sector stratum, zero on
 \textsc{uniform-weight}.
\item\label{td-iii} \emph{Symmetry.}
 $K_{w_1, w_2} = K_{w_2, w_1}$ (bosonic sector) or
 $(-1)^{|w_1||w_2|} K_{w_2, w_1}$ (super-graded sector), so
 $K(\mathcal{A})$ lands in $\operatorname{Sym}^2$ or
 $\operatorname{Sym}^2_{\mathrm{super}}$ accordingly.
\item\label{td-iv} \emph{Tensor chiral Mumford formula.}
$
 \int_{\overline{\mathcal{M}}_{g,n}} \operatorname{tr}_\mathcal{F}
 \bigl(\exp(K(\mathcal{A}))\bigr)
 \;=\;
 \prod_w
 \int_{\overline{\mathcal{M}}_{g,n}} \exp(\kappa_w \cdot
 \omega_g^{(w)})
 \;+\; \delta^{\mathrm{cross}}_g,
$
\noindent
generalising Mumford's scalar formula; the first factor is the
weight-diagonal contribution, $\delta^{\mathrm{cross}}_g$ the
cross-channel correction from the off-diagonal entries.
\end{enumerate}
\end{theorem}

\begin{proof}[Sketch]
\textit{Step 1 (tensor $d^2$ computation).} Write the modular bar
differential $d$ in the weight-graded decomposition
$d = \sum_{w_1, w_2} d_{w_1 \to w_2}$ where $d_{w_1 \to w_2}$
shifts weight from $w_1$ to $w_2$ via a chiral contraction
residue. Then
$d^2 = \sum_{w_1, w_3} \bigl(
 \sum_{w_2} d_{w_2 \to w_3} d_{w_1 \to w_2}\bigr)
 = \sum_{w_1, w_3} K_{w_1, w_3}$
lives naturally in the weight-pair grading
$\mathcal{F}_{w_1}^\vee \otimes \mathcal{F}_{w_3}^\vee$.

\textit{Step 2 (diagonal entry = UNIFORM-WEIGHT scalar).} When
$w_1 = w_3 = w$, the contracting residues preserve weight and
produce a scalar multiple of the weight-$w$ Arakelov form
$\omega_g^{(w)}$: this is the weight-$w$ restriction of Vol~I's
scalar Mumford--Arakelov formula $d^2 = \kappa \cdot \omega_g$:
at genus \(1\) this is the twisted tensor criterion, while the
higher-genus continuation uses the conditional curved-Dunn
\(H^2\) bridge
(Theorem~\ref{thm:curved-dunn-H2-vanishing-all-genera}).

\textit{Step 3 (off-diagonal entry = cross-channel).} When
$w_1 \ne w_3$, the residues combining $d_{w_1 \to w_2}$ and
$d_{w_2 \to w_3}$ produce a cross-channel $2$-form identified with
$\delta F_g^{\mathrm{cross}}|_{w_1, w_3}$ by direct comparison
with the multi-weight GRR-type anomaly computation of Vol~I, in the
all-weight scope where every $\mathrm{obs}_g$, $F_g$, and
$\lambda_g$ carries its explicit weight assignment.

\textit{Step 4 (symmetry).} The graded commutativity of the
bar differential's tensor components follows from the
Getzler--Kapranov modular operad axioms: vertical composition of
residues commutes with horizontal composition up to a Koszul
sign, yielding $K_{w_1, w_2} = (-1)^{|w_1||w_2|} K_{w_2, w_1}$.

\textit{Step 5 (tensor Mumford formula).} Pulling back the
characteristic classes $c_1(\det \mathcal{F}_w)$ along the
clutching morphisms gives a factorisation of the total class
$\exp(K)$ into weight-diagonal factors
$\exp(\kappa_w \omega_g^{(w)})$ plus cross-channel corrections.
Integrating against the fundamental class of
$\overline{\mathcal{M}}_{g,n}$ yields the tensor chiral Mumford
formula.
\end{proof}

\begin{example}[Virasoro: diagonal tensor curvature for
 UNIFORM-WEIGHT = scalar]
\label{ex:theoremD-virasoro}
For $\mathcal{A} = \mathrm{Vir}_c$ with $c \neq 0$, the weight
bundle is $\mathcal{F} = \bigoplus_{w \ge 0} \mathrm{Vir}_c^{(w)}$
with $\mathrm{Vir}_c^{(w)}$ the conformal-block local system at
weight $w$. Virasoro has a single strong generator $T$ at
weight~$2$, so $\kappa_2(\mathrm{Vir}_c) = c/2$ and
$\kappa_w(\mathrm{Vir}_c) = 0$ for all $w \neq 2$ (single-generator
concentration, NOT a linear-in-weight law).
On UNIFORM-WEIGHT locus: only the $w_1 = w_2 = 2$ diagonal entry
contributes; total trace
$\operatorname{tr}_\mathcal{F}(K) = \kappa_2 \cdot \omega_g^{(2)} =
(c/2) \cdot \omega_g^{(2)}$ recovers Vol~I's scalar Theorem~D. On ALL-WEIGHT stratum: off-diagonal
entries $K_{w_1, w_2}$ (with $w_1 + w_2 = 4$, e.g.\ the quartic
contact $\mathfrak{Q}^{\mathrm{contact}}_{\mathrm{Vir}} =
10/[c(5c+22)]$) enter at genus $g \ge 2$; the tensor Mumford
formula reproduces the shadow tower's quartic coefficient.
\end{example}

\begin{remark}[Scope and Vol~I comparison]
\label{rem:theoremD-tensor-scope}
Theorem~\ref{thm:theoremD-tensor-arakelov} holds unconditionally
on the \textsc{uniform-weight} stratum, diagonal entries
weight-by-weight reducing to the scalar Vol~I formula. On the
\textsc{all-weight} stratum with cross-channel weight couplings,
off-diagonal $K_{w_1, w_2}$ requires the curved-Dunn $H^2$ bridge
(Theorem~\ref{thm:curved-dunn-H2-vanishing-all-genera}); with the
modular-bootstrap acyclicity and degree-\(2\) comparison
hypotheses, the tensor formula extends to all genera for
weight-bounded chiral algebras in classes
$\mathbf{G}, \mathbf{L}, \mathbf{C}, \mathbf{M}$. The
Feigin--Frenkel critical-level stratum ($\kappa = 0$) degenerates
to $K = 0$, with the anomaly carried by the Feigin--Frenkel-centre
cocycle (class~$\mathbf{FF}$).
\end{remark}

\begin{remark}[Identification of the weight-bundle ambient
 $\mathcal{F}$]
\label{rem:theoremD-ambient-identification}
Three ambients carry tensor targets
$\operatorname{Sym}^2(F^\vee) \otimes \Omega^2$:
(A) the Hodge bundle $\Lambda_g$ of rank $g$;
(B) the $\Lambda$-twisted tensor ladder
$\bigotimes_k \Lambda_g^{\otimes k}$, an infinite sum without
finite-rank structure;
(C) the weight-bundle
$\mathcal{F} = \bigoplus_w \mathcal{F}_w$. Only (C) carries the
rank structure $K$ demands: (A) of rank $g$ is too coarse for the
weight separation visible to the bar differential; (B) has no
finite rank at each degree. Arakawa $C_2$-cofiniteness gives
finite-dim $\mathcal{F}_w$ at each $w$; the Li-filtration
weight-truncation makes $\bigoplus_w \mathcal{F}_w$ a graded local
system of finite total rank over each stable stratum.
\end{remark}

\begin{proposition}[Family-specific weight distribution of the
 diagonal trace]
\label{prop:theoremD-weight-distribution}
\ClaimStatusProvedHere
The identification
$\operatorname{tr}_\mathcal{F}(K) = \kappaChHodge(\cA) \cdot \omega_g$
holds across the Vol~I landscape; the per-weight decomposition
$\kappa_w$ is family-specific:
\begin{itemize}
\item Heisenberg $\mathcal{H}_k$ (single strong generator $J$ at
 weight~$1$): $\kappa_1 = k$, all other $\kappa_w = 0$;
 $\operatorname{tr}(K) = k = \kappaChHodge(\mathcal{H}_k)$.
\item Virasoro $\mathrm{Vir}_c$ (single strong generator $T$ at
 weight~$2$): $\kappa_2 = c/2$, all other $\kappa_w = 0$;
 $\operatorname{tr}(K) = c/2 = \kappaChHodge(\mathrm{Vir}_c)$.
\item $\mathcal{W}_N$ (strong generators at weights
 $2, 3, \ldots, N$): $\kappa_j = c/j$ for
 $j \in \{2, \ldots, N\}$, all other $\kappa_w = 0$;
 $\operatorname{tr}(K) = \sum_{j=2}^N c/j = c(H_N - 1) =
 \kappaChHodge(\mathcal{W}_N)$.
\end{itemize}
Virasoro refutes a uniform linear-in-weight law: a linear
distribution would predict $\kappa_2 = c$, in conflict with
$\kappaChHodge(\mathrm{Vir}_c) = c/2$.
\end{proposition}

\begin{proof}[Proof]
Heisenberg and Virasoro each have a single strong generator;
single-weight support follows from the chiral bar differential
vanishing on weight cohomology outside the generator's weight. The
principal $\mathcal{W}_N$ generates strong conformal primaries at
weights $2, 3, \ldots, N$ (Fateev--Lukyanov), each contributing
$c/j$ through the Zamolodchikov OPE at weight $j$; the harmonic
sum $\sum_{j=2}^N c/j$ recovers the Vol~I kappa formula
$\kappaChHodge(\mathcal{W}_N) = c(H_N - 1)$ with
$H_N = \sum_{j=1}^N 1/j$ (C4 of the Vol~I census). Substituting
$\kappa_w = \alpha w$ into Virasoro gives $\kappa_2 = 2\alpha$;
solving against $\kappaChHodge(\mathrm{Vir}_c) = c/2$ yields
$\alpha = c/4$, whence $\kappa_3 = 3c/4 \neq 0$, in conflict with
single-generator support.
\end{proof}

\section{Corollaries and cross-volume coherence}
\label{sec:theorems-C-D-cross-volume}

The Vol~II-native total forms contain Vol~I's genus-wise statements
as partial-to-total restrictions:

\begin{itemize}
 \item Vol~I Theorem~\ref*{thm:theorem-C} (genus-wise Lagrangian)
 restricts Theorem~C to each stable stratum of
 $\overline{\cM}_{g,n}$; the total Lagrangian glues genus-wise
 sections along clutching morphisms
 (Theorem~\ref{thm:theoremC-total-shifted-symplectic}(iii)).
 \item Vol~I Theorem~\ref*{thm:theorem-D} (uniform-weight scalar
 $\obs_g = \kappaChHodge(\cA) \cdot \lambda_g$) is the total
 diagonal trace
 $\operatorname{tr}_\cF(K) = \sum_w \kappa_w(\cA) \cdot
 \omega_g^{(w)}$ collapsing to
 $\kappaChHodge(\cA) \cdot \lambda_g$ by the family-specific
 distribution law
 (Proposition~\ref{prop:theoremD-weight-distribution}); the
 off-diagonal $\delta F_g^{\mathrm{cross}}$ carries the
 cross-channel couplings at $w_1 \neq w_2$, computed for $W_3$ at
 $g = 2$ in Chapter~\ref{ch:theory-hochschild}.
\end{itemize}

Section~C inscribes the total $-(3g{-}3{+}n)$-shifted symplectic
structure on the characteristic bundle; Section~D inscribes the
tensor curvature on the weight-bundle. Both close in clutching
compatibility and recover the Vol~I scalar form by specialisation.
The characteristic-datum package thus closes at the tensor level,
realising the genus-wise relations of Vol~I as global sections of
shifted-symplectic and tensor-valued ambients.
